% id: 41-02
% book: 베이직쎈 미적분1 · 02 함수의 연속
% section: 학교 시험 기출로 실전 감각 UP!
% figure: tikz:fig-41-02
% answer: 
% answer_source: 
% variant_level: 0 (원본 전사)
% 이 파일은 build.mjs 가 items.json 에서 만든다. 고칠 때는 items.json 을 고친다.
\smash{\raisebox{1.5mm}{\dmkichul{베이직쎈 미적분1 41쪽 02번 · 꼭 나와요}}}
\noindent\begin{minipage}[t]{0.58\linewidth}\vspace{0pt}\raggedright 함수 $y=f(x)$의 그래프가 오른쪽 그림과 같을 때, 다음 중 옳지 \underline{않은} 것은?\end{minipage}\hfill\begin{minipage}[t]{0.4\linewidth}\vspace{0pt}\centering \resizebox{\ifdim\width>\linewidth\linewidth\else\width\fi}{!}{% 41-02(41쪽) — x=1·x=2 에서 끊어진 y=f(x) 의 그래프. 스캔 화질이 낮아 TikZ 로 다시 그림.
% 원문에 있는 것만: 채운 점 (1,1)·(2,0) · 빈 점 (1,-1)·(2,1) · 안내 점선 여섯 줄 · 라벨 -1, O, 1, 2, 3, 1, -1, x, y, y=f(x).
\begin{tikzpicture}[x=0.60cm, y=0.56cm, line width=0.5pt]
  \draw[->, >=stealth, line width=0.7pt] (-2.2,0) -- (4.1,0) node[below=1pt, font=\small] {$x$};
  \draw[->, >=stealth, line width=0.7pt] (0,-1.75) -- (0,2.45) node[left=1pt, font=\small] {$y$};
  \draw[densely dotted, line width=0.4pt] (-1,0) -- (-1,1);
  \draw[densely dotted, line width=0.4pt] (-1,1) -- (2,1);
  \draw[densely dotted, line width=0.4pt] (1,1) -- (1,-1);
  \draw[densely dotted, line width=0.4pt] (0,-1) -- (1,-1);
  \draw[densely dotted, line width=0.4pt] (2,1) -- (2,0);
  \draw[densely dotted, line width=0.4pt] (3,1) -- (3,0);
  \draw[line width=0.6pt] (-2.05,2.05) -- (1,-1);
  \draw[line width=0.6pt] (1,-1) -- (2,0);
  \draw[line width=0.6pt] (2,1) -- (3.85,1);
  \fill (1,1) circle (2.2pt);
  \fill (2,0) circle (2.2pt);
  \draw[fill=white, line width=0.5pt] (1,-1) circle (2.2pt);
  \draw[fill=white, line width=0.5pt] (2,1) circle (2.2pt);
  \node[left=1pt, font=\small] at (0,1.22) {$1$};
  \node[left=1pt, font=\small] at (0,-1) {$-1$};
  \node[below=1pt, font=\small] at (-1,0) {$-1$};
  \node[below left=-2pt and -3pt, font=\small] at (0,0) {$\pt{O}$};
  \node[below=1pt, font=\small] at (1,0) {$1$};
  \node[below=1pt, font=\small] at (2,0) {$2$};
  \node[below=1pt, font=\small] at (3,0) {$3$};
  \node[font=\small] at (1.95,1.6) {$y=f(x)$};
\end{tikzpicture}
}\end{minipage}\par
\par\medskip\par\noindent\hangindent=1.4em\hangafter=1 {\small ①}\ $\displaystyle\lim_{x\to 1} f(x)$의 값이 존재한다.\par\noindent\hangindent=1.4em\hangafter=1 {\small ②}\ $\displaystyle\lim_{x\to 2} f(x)$의 값이 존재하지 않는다.\par\noindent\hangindent=1.4em\hangafter=1 {\small ③}\ 함수 $f(x)$는 $x=1$에서 불연속이다.\par\noindent\hangindent=1.4em\hangafter=1 {\small ④}\ 함수 $f(x)$는 $x=3$에서 연속이다.\par\noindent\hangindent=1.4em\hangafter=1 {\small ⑤}\ $-1\le x\le 3$에서 함수 $f(x)$가 불연속인 $x$의 값은 1개이다.\par\medskip
